\section*{第 \ref{ch:b1:diffeq} 章 --- 线性微分方程}

\begin{solution}{exo:b1:diffeq:1}
齐次解：$y_h = \lambda\,\eu^{-2x}$。特解：试取 $y_p =
c\,\eu^{3x}$（$3$ 不是 $r + 2$ 的根）：$3c + 2c = 1$，$c =
\frac15$。通解：$y = \frac{\eu^{3x}}{5} +
\lambda\,\eu^{-2x}$。由 $y(0) = 1$ 得 $\frac15 + \lambda = 1$，
$\lambda = \frac45$，于是 $y = \frac{\eu^{3x} + 4\,\eu^{-2x}}{5}$。
\end{solution}

\begin{solution}{exo:b1:diffeq:2}
在 $\intoo{0}{+\infty}$ 上，两边除以 $x$：$y' - \frac{2}{x}\,y = x^2$。
此处 $A(x) = -2\ln x$，$\eu^{-A(x)} = x^2$：齐次解为
$\lambda x^2$。由常数变易法，$\mu'(x) = x^2 \cdot x^{-2} =
1$，故 $\mu = x + \lambda$，于是
\[
y(x) = x^3 + \lambda x^2, \qquad \lambda \in \R .
\]
\emph{验证：}$x(3x^2 + 2\lambda x) - 2(x^3 + \lambda x^2) = x^3$。
\end{solution}

\begin{solution}{exo:b1:diffeq:3}
$y'' - 3y' + 2y = 0$：根为 $1$ 与 $2$；$\;y = \lambda\,\eu^{x} +
\mu\,\eu^{2x}$。

$y'' + 4y' + 4y = 0$：二重根 $-2$；$\;y = (\lambda + \mu
x)\,\eu^{-2x}$。

$y'' - 2y' + 5y = 0$：根为 $1 \pm 2\iu$；$\;y = \eu^{x}(\lambda\cos
2x + \mu\sin 2x)$。
\end{solution}

\begin{solution}{exo:b1:diffeq:4}
齐次解：根为 $\pm 1$，$y_h = \lambda\,\eu^x + \mu\,\eu^{-x}$。
右端为多项式（$\gamma = 0$ 不是根），设特解：
$y_p = ax^2 + bx + c$；代入得 $2a - (ax^2 + bx + c) = x^2$，
比较系数得 $a = -1$，$b = 0$，$c = 2a = -2$，即 $y_p = -x^2 - 2$。通解为
$y = -x^2 - 2 + \lambda\eu^x + \mu\eu^{-x}$。

柯西问题：$y(0) = -2 + \lambda + \mu = 0$，且 $y'(0) = \lambda - \mu =
1$，得 $\lambda = \frac32$，$\mu = \frac12$。于是
$y = -x^2 - 2 + \frac{3\eu^x + \eu^{-x}}{2}$。
\end{solution}

\begin{solution}{exo:b1:diffeq:5}
$a(x) = \tan x$，$A(x) = -\ln(\cos x)$（该区间上 $\cos > 0$，
故有意义），$\eu^{-A} = \cos x$：齐次解为 $\lambda\cos x$。
由常数变易法，$\mu'(x) = \sin 2x \cdot \frac{1}{\cos x} =
2\sin x$，故 $\mu = -2\cos x + \lambda$，于是
\[
y(x) = -2\cos^2 x + \lambda \cos x .
\]
\emph{验证：}$y' = 4\cos x \sin x - \lambda\sin x$，且 $y\tan x =
-2\cos x\sin x + \lambda \sin x$，二者之和为 $2\cos x\sin x =
\sin 2x$，与要求相符。
\end{solution}

\begin{solution}{exo:b1:diffeq:6}
$\chi(r) = r^2 - 4r + 3 = (r-1)(r-3)$：$\gamma = 1$ 为单根
（$m = 1$）。试取 $y_p = x(ax + b)\,\eu^x$。记 $u = ax^2 + bx$，
\[
y_p'' - 4y_p' + 3y_p = \bigl(u'' + (2 - 4)u' + \chi(1) u\bigr)\eu^x
= \bigl(2a - 2(2ax + b)\bigr)\eu^x .
\]
与 $(2x + 1)\eu^x$ 比较系数：$-4a = 2$，$2a - 2b = 1$，于是 $a =
-\frac12$，$b = -1$。通解为
\[
y = -\Bigl(\frac{x^2}{2} + x\Bigr)\eu^{x} + \lambda\,\eu^{x} +
\mu\,\eu^{3x}, \qquad (\lambda, \mu) \in \R^2 .
\]
\end{solution}

\begin{solution}{exo:b1:diffeq:7}
齐次解：$y_h = \lambda\cos 2x + \mu\sin 2x$。

右端为 $x$（$\gamma = 0$ 不是根），设 $y_1 = ax + b$，
代入得 $4(ax + b) = x$，即 $y_1 = \frac x4$。

右端为 $\sin 2x = \Im(\eu^{2\iu x})$，且 $2\iu$ 是
$r^2 + 4$ 的单根：试取 $z = c\,x\,\eu^{2\iu x}$，则 $z'' + 4z =
4\iu c\,\eu^{2\iu x}$，令其等于 $\eu^{2\iu x}$，取 $c =
\frac{1}{4\iu} = -\frac{\iu}{4}$，于是 $z = -\frac{\iu x}{4}(\cos 2x +
\iu \sin 2x)$，故 $y_2 = \Im(z) = -\frac{x\cos 2x}{4}$。

由叠加原理，
\[
y = \frac{x}{4} - \frac{x\cos 2x}{4} + \lambda\cos 2x + \mu\sin 2x .
\]
\end{solution}

\begin{solution}{exo:b1:diffeq:8}
方程 $T' + kT = 20k$ 的常值特解为 $20$，
齐次解为 $\lambda\eu^{-kt}$，故 $T(t) = 20 +
\lambda\,\eu^{-kt}$。由 $T(0) = 80$ 得 $\lambda = 60$：
\[
T(t) = 20 + 60\,\eu^{-kt} .
\]
由 $T(10) = 50$ 得 $\eu^{-10k} = \frac12$，故 $k = \frac{\ln 2}{10}$。于是
$T(t) = 25$ 要求 $\eu^{-kt} = \frac{5}{60} = \frac{1}{12}$，即
\[
t = \frac{\ln 12}{k} = 10\,\frac{\ln 12}{\ln 2} \approx 35.8
\text{ 分钟.}
\]
\end{solution}

\begin{solution}{exo:b1:diffeq:9}
\begin{enumerate}
  \item $\chi(r) = r^2 + 2\varepsilon r + 1$，$\Delta = 4(\varepsilon^2
        - 1)$。
        当 $\varepsilon \in \intco{0}{1}$ 时，根为 $-\varepsilon \pm
        \iu\sqrt{1 - \varepsilon^2}$，故
        $y = \eu^{-\varepsilon t}\bigl(\lambda\cos\omega t +
        \mu\sin\omega t\bigr)$，其中 $\omega = \sqrt{1 -
        \varepsilon^2}$。
        当 $\varepsilon = 1$ 时，二重根为 $-1$，$y = (\lambda + \mu
        t)\,\eu^{-t}$。
        当 $\varepsilon > 1$ 时，实根为 $r_\pm = -\varepsilon \pm
        \sqrt{\varepsilon^2 - 1}$，二者均 $< 0$，且 $y =
        \lambda\eu^{r_+t} + \mu\eu^{r_-t}$。
  \item 当 $\varepsilon \in \intoo{0}{1}$ 时，$\abs y \leq
        \eu^{-\varepsilon t}(\abs\lambda + \abs\mu) \to 0$。当
        $\varepsilon = 1$ 时，$(\lambda + \mu t)\eu^{-t} \to 0$
        （指数因子压过多项式，
        见\cref{prop:b1:functions:powerrules}）。当 $\varepsilon > 1$ 时，
        由 $r_\pm < 0$ 知两个指数项均衰减（事实上
        $\sqrt{\varepsilon^2 - 1} < \varepsilon$）。当 $\varepsilon =
        0$ 时，$y = \lambda\cos t + \mu\sin t$ 的振幅恒为
        $\sqrt{\lambda^2 + \mu^2} \neq 0$，除非 $y = 0$。
  \item 写成 $\lambda\cos\omega t + \mu\sin\omega t =
        R\cos(\omega t - \varphi)$，其中 $R = \sqrt{\lambda^2 + \mu^2}
        > 0$。$y$ 的零点与 $\cos(\omega t -
        \varphi)$ 的零点相同（因子 $\eu^{-\varepsilon t}$ 永不
        为零）：$\omega t - \varphi \equiv \frac\pi2 \pmod \pi$，
        即零点构成公差为 $\frac{\pi}{\omega} =
        \frac{\pi}{\sqrt{1 - \varepsilon^2}}$ 的等差数列。
\end{enumerate}
\end{solution}

\begin{solution}{exo:b1:diffeq:10}
设 $x = y = 0$：$2f(0) = 2f(0)^2$，由 $f(0) \neq 0$ 得 $f(0) =
1$。固定 $x$，对方程关于 $y$ 求两次导数：
\[
f''(x+y) + f''(x-y) = 2 f(x) f''(y) .
\]
令 $y = 0$：$\;2f''(x) = 2 f(x) f''(0)$，即
\[
f''(x) = c\,f(x), \qquad c = f''(0).
\]
\emph{情形 $c = \omega^2 > 0$：}$f(x) = \lambda\cosh\omega x +
\mu\sinh\omega x$；由 $f(0) = 1$ 得 $\lambda = 1$。代入
函数方程，并利用加法公式
（见\cref{prop:b1:functions:hyprules}），方程迫使 $\mu = 0$
（比较 $\sinh\omega x \sinh\omega y$ 的系数，或
取 $x = y$ 代入）：$f = \cosh\omega x$，它确实满足
$\cosh(x+y) + \cosh(x-y) = 2\cosh x\cosh y$。

\emph{情形 $c = -\omega^2 < 0$：}类似地，$f(x) = \cos\omega x$
（$\omega \neq 0$），它满足该方程。

\emph{情形 $c = 0$：}$f$ 为仿射函数，且 $f(0) = 1$，故 $f(x) = 1 + \mu x$；
方程迫使 $\mu = 0$，故舍去（$f$ 不为常数）。

结论：解为 $f(x) = \cos\omega x$ 与 $f(x) =
\cosh\omega x$，其中 $\omega > 0$。
\end{solution}

\begin{solution}{exo:b1:diffeq:11}
令 $z(t) = y(\eu^t)$，于是 $y(x) = z(\ln x)$，其中 $x > 0$。则
\[
y'(x) = \frac{z'(\ln x)}x,
\qquad
y''(x) = \frac{z''(\ln x) - z'(\ln x)}{x^2} ,
\]
代入方程，得
\[
x^2y'' - xy' + y = \bigl(z'' - z'\bigr) - z' + z
= z'' - 2z' + z = 0 .
\]
特征多项式为 $(r - 1)^2$，二重根为 $1$，故 $z(t) =
(\lambda + \mu t)\,\eu^t$，再回到变量 $x = \eu^t$，得
\[
y(x) = (\lambda + \mu\ln x)\,x,
\qquad \lambda, \mu \in \R .
\]
\end{solution}

\begin{solution}{exo:b1:diffeq:12}
\begin{enumerate}
  \item 化为标准形式 $y' - \frac2x\,y = 0$（在每个区间上）：
        $A(x) = -2\ln\abs x$，故解为 $y = a x^2$，定义于
        $\intoo0{+\infty}$ 上，及 $y = b x^2$，定义于 $\intoo{-\infty}0$ 上，
        两个常数相互独立
        （见\cref{thm:b1:diffeq:homogeneous1}）。
  \item 设 $y = ax^2$（$x \geq 0$）、$bx^2$（$x < 0$）。
        在每个开半直线上 $y$ 可导，且 $xy' = 2y$。
        在 $0$ 处：差商 $\frac{y(h) - y(0)}h =
        ah$ 或 $bh$ 趋于 $0$，故 $y'(0) = 0$ 存在，且
        方程在 $x = 0$ 处化为 $0 \cdot y'(0) = 2y(0) = 0$：
        成立。故 $y$ 在整个 $\R$ 上满足方程。
  \item $\R$ 上的解恰为这些拼接而成的函数：对一阶方程而言，它们
        构成一个\emph{两}参数族。
        这与\cref{thm:b1:diffeq:voc}并无矛盾，
        该定理的条件在此不成立：写成 $y' + a(x)y = 0$ 时，
        系数 $a(x) = -\frac2x$ 在 $0$ 处不连续，
        --- 事实上甚至无定义 --- 故 $\R$ 不是
        该定理适用的区间。奇点 $0$ 将两条半直线
        隔开，$y(0) = 0$ 是被迫取得的值，
        两侧之间不传递任何信息。每个柯西
        数据（在 $x_0 \neq 0$ 处给定）只决定
        其所在半直线上的解，即包含 $x_0$ 的半直线。
\end{enumerate}
\end{solution}

\begin{solution}{pb:b1:diffeq:1}
\textbf{1.} $\chi(r) = r^2 + 2\lambda r + \omega_0^2$，$\Delta =
4(\lambda^2 - \omega_0^2) < 0$，当 $0 < \lambda < \omega_0$ 时，根为
$-\lambda \pm \iu\omega_d$，由
\cref{thm:b1:diffeq:homogeneous2}可知
\[
x(t) = \eu^{-\lambda t}\bigl(\lambda_1\cos\omega_d t +
\mu_1\sin\omega_d t\bigr), \qquad (\lambda_1, \mu_1) \in \R^2 .
\]
当 $\lambda = 0$ 时：$x = \lambda_1\cos\omega_0 t +
\mu_1\sin\omega_0 t$。临界情形（$\lambda = \omega_0$）和
过阻尼情形（$\lambda > \omega_0$）的解
与\cref{exo:b1:diffeq:9}中的解相同（经时间尺度变换后）：$(\lambda_1 +
\mu_1 t)\eu^{-\lambda t}$，分别为 $\eu^{r_\pm t}$ 的线性组合，
其中 $r_\pm = -\lambda \pm \sqrt{\lambda^2 - \omega_0^2}$。

\textbf{2.} 欠阻尼情形：$\abs x \leq \eu^{-\lambda
t}(\abs{\lambda_1} + \abs{\mu_1}) \to 0$。临界情形：$(\lambda_1 +
\mu_1 t)\eu^{-\lambda t} \to 0$，因为指数因子胜过
多项式（见\cref{prop:b1:functions:powerrules}）。过阻尼情形：
$r_- < r_+ = -\lambda + \sqrt{\lambda^2 - \omega_0^2} < 0$
因为 $\sqrt{\lambda^2 - \omega_0^2} < \lambda$，两个
指数项均衰减。

\textbf{3.} 沿 $(H)$ 的解，利用 $x'' = -2\lambda x' -
\omega_0^2 x$，有
\[
\mathcal E'(t) = x'x'' + \omega_0^2 x x'
= x'\bigl(-2\lambda x' - \omega_0^2 x\bigr) + \omega_0^2 xx'
= -2\lambda\,x'^2 \leq 0 .
\]
若 $x(t_0) = x'(t_0) = 0$，则 $\mathcal E(t_0) = 0$；$\mathcal E$
非负且不增，故 $\mathcal E \equiv 0$ 于
$\intco{t_0}{+\infty}$ 上，从而 $x \equiv 0$ 于该区间；当 $t \leq
t_0$ 时，对 $\tilde x(t) = x(2t_0 - t)$ 应用同一论证，它
满足阻尼系数为 $-\lambda$ 的方程，但仍有
$\tilde{\mathcal E}(t_0) = 0$，且 $\tilde{\mathcal E}' =
+2\lambda\tilde x'^2 \geq 0$，并有 $\tilde{\mathcal E} \geq 0$；
非负、不减，且在
$\intoc{-\infty}{t_0}$ 的右端为零，故整段为零。于是 $x \equiv 0$ 在
$\R$ 上 --- 这是唯一性的一个能量证明，对一切 $\lambda
\geq 0$ 均有效。

\textbf{4.} $x(t + T_d) = R\,\eu^{-\lambda t}\eu^{-\lambda
T_d}\cos(\omega_d t + 2\pi - \varphi) = \eu^{-\lambda T_d}x(t)$.
每个伪周期的缩减因子为 $\eu^{-\delta}$，
其中 $\delta = \lambda T_d = \frac{2\pi\lambda}{\omega_d}$。对于
$\omega_0 = 1$，$\lambda = 0.1$：$\omega_d = \sqrt{0.99} =
0.99499$，故 $\delta = \frac{0.62832}{0.99499} = 0.6315$：每
次摆动保留其振幅的 $\eu^{-0.63} \approx 53\%$。

\textbf{5.} 振幅因子为 $\eu^{-\lambda t}$，衰减 $\eu$ 倍所需的
时间为 $t = \frac1\lambda$。该区间包含
$\frac{1/\lambda}{T_d} = \frac{\omega_d}{2\pi\lambda}$ 个伪周期。
当 $\lambda \ll \omega_0$ 时，$\omega_d \approx
\omega_0$，其值 $\approx \frac{\omega_0}{2\pi\lambda} =
\frac Q\pi$。$Q = 300$ 的吉他弦大约振动一百个
周期；$Q = 1$ 的门阻尼器则连一个周期都无法完成。

\textbf{6.} 将 $z\,\eu^{\iu\Omega t}$ 代入左端，
得 $z\,(-\Omega^2 + 2\iu\lambda\Omega +
\omega_0^2)\,\eu^{\iu\Omega t}$，它恰好等于 $A\,\eu^{\iu\Omega
t}$，当且仅当 $z = \frac{A}{\omega_0^2 - \Omega^2 +
2\iu\lambda\Omega}$（分母非零，其虚部
为 $2\lambda\Omega > 0$）。由于系数为实数，
取实部 $x_p = \Re\bigl(z\eu^{\iu\Omega t}\bigr)$，它是
右端为 $\Re\bigl(A\eu^{\iu\Omega t}\bigr) =
A\cos\Omega t$ 的方程的解。

\textbf{7.} 记 $\omega_0^2 - \Omega^2 + 2\iu\lambda\Omega =
\sqrt{D}\,\eu^{\iu\varphi}$，其中 $D = (\omega_0^2 - \Omega^2)^2 +
4\lambda^2\Omega^2$，且 $\varphi \in \intoo0\pi$
（虚部 $2\lambda\Omega$ 为正），从而 $\tan\varphi =
\frac{2\lambda\Omega}{\omega_0^2 - \Omega^2}$。于是 $z =
\frac{A}{\sqrt D}\eu^{-\iu\varphi}$，且
\[
x_p(t) = \Re\Bigl(\frac A{\sqrt D}\,\eu^{\iu(\Omega t -
\varphi)}\Bigr) = R\cos(\Omega t - \varphi),
\qquad R = \frac A{\sqrt D} .
\]

\textbf{8.} 当 $\Omega \to 0^+$ 时：$D \to \omega_0^4$，故 $R \to
A/\omega_0^2$，且 $\tan\varphi \to 0^+$，其中 $\varphi \in
\intoo0{\frac\pi2}$：$\varphi \to 0$。质量随外力
准静态地运动，位移为外力除以劲度系数。当 $\Omega \to
+\infty$：$D \sim \Omega^4$，故 $R \sim A/\Omega^2 \to 0$，且
$\varphi \to \pi$（复数 $\omega_0^2 - \Omega^2 +
2\iu\lambda\Omega$ 落在第二象限，辐角 $\to
\pi$）：质量几乎不动，且与外力反相——
惯性占主导。

\textbf{9.} 由 \cref{thm:b1:diffeq:structure2}~(1) 可知，每个解
均为 $x = x_p + x_h$，其中 $x_h$ 满足 $(H)$；由第 2 题，$x_h(t)
\to 0$，故 $x(t) - x_p(t) \to 0$：所有解趋于
同一稳态。初值只影响瞬态部分。

\textbf{10.} 若 $x$ 为周期解，则 $x - x_p = x_h$ 为
$(H)$ 的周期解，且以 $0$ 为其在 $+\infty$ 处的极限；一个以
$0$ 为极限的周期函数恒等于 $0$（其在一个周期上的值
无限重复，
故每个值都是某个趋于 $0$ 的子列的极限。因此 $x = x_p$。

\textbf{11.} 此时 $\lambda = 1$，$\omega_0^2 = 2$，$\Omega = 1$，
$A = 1$：$z = \frac1{2 - 1 + 2\iu} = \frac{1 - 2\iu}5$，于是
\[
x_p = \Re\Bigl(\frac{(1 - 2\iu)(\cos t + \iu\sin t)}5\Bigr)
= \frac{\cos t + 2\sin t}5 .
\]
齐次部分：$r^2 + 2r + 2$ 的根为 $-1 \pm \iu$，于是 $x_h =
\eu^{-t}(C\cos t + S\sin t)$。由初始条件 $x(0) = \frac15 + C = 0$
得 $C = -\frac15$；求导得 $x'(0) = \frac25 - C + S =
0$，得 $S = C - \frac25 = -\frac35$。于是
\[
x(t) = \underbrace{\frac{\cos t + 2\sin t}5}_{\text{稳态}}
- \underbrace{\eu^{-t}\,\frac{\cos t + 3\sin
t}5}_{\text{瞬态}} ,
\]
瞬态项随 $\eu^{-t}$ 衰减。

\textbf{12.} 展开得 $g(u) = u^2 - 2(\omega_0^2 - 2\lambda^2)u
+ \omega_0^4 = (u - u_r)^2 + g(u_r)$，其中 $u_r = \omega_0^2 -
2\lambda^2$，且
\[
g(u_r) = \omega_0^4 - u_r^2 = (\omega_0^2 - u_r)(\omega_0^2 +
u_r) = 2\lambda^2\bigl(2\omega_0^2 - 2\lambda^2\bigr) =
4\lambda^2(\omega_0^2 - \lambda^2) .
\]
若 $2\lambda^2 < \omega_0^2$，则 $u_r > 0$ 为允许的
平方频率：$g$ 在该点取严格极小值，故 $R = A/\sqrt
g$ 在 $\Omega_r = \sqrt{u_r} =
\sqrt{\omega_0^2 - 2\lambda^2}$ 处取严格极大值，且 $R_{\max} = A/\sqrt{g(u_r)}
= \frac{A}{2\lambda\sqrt{\omega_0^2 - \lambda^2}}$。

\textbf{13.} $R(0) = A/\omega_0^2$，故
\[
\frac{R_{\max}}{R(0)}
= \frac{\omega_0^2}{2\lambda\sqrt{\omega_0^2 - \lambda^2}}
= \frac{\omega_0}{2\lambda}\cdot
\frac{\omega_0}{\sqrt{\omega_0^2 - \lambda^2}}
= Q\,\Bigl(1 - \frac{\lambda^2}{\omega_0^2}\Bigr)^{-1/2}
\geq Q .
\]
弱阻尼时修正因子接近 $1$：共振
使静态位移基本放大为 $Q$ 倍。

\textbf{14.} $V(\Omega)^2 = \frac{A^2 u}{g(u)}$，其中 $u =
\Omega^2$。其导数的符号与 $g(u) - u\,g'(u) =
(\omega_0^2 - u)^2 + 4\lambda^2 u - u\bigl(2(u - \omega_0^2) +
4\lambda^2\bigr) = (\omega_0^2 - u)^2 + 2u(\omega_0^2 - u) =
(\omega_0^2 - u)(\omega_0^2 + u)$，当 $u < \omega_0^2$ 时为正，
超出后为负，故严格极大值恰在 $\Omega =
\omega_0$ 处取得，对任意阻尼均成立。此时 $\tan\varphi$ 趋于无穷，
在 $\varphi \in \intoo0\pi$ 中，$\varphi = \frac\pi2$，且 $x_p'(t) =
-R\Omega\sin(\Omega t - \frac\pi2) = R\Omega\cos(\Omega t)$，与外力
恰同相，此时功率传递最优。

\textbf{15.} $R = R_{\max}/\sqrt2 \iff g(u) = 2g(u_r) \iff (u -
u_r)^2 = g(u_r) = 4\lambda^2(\omega_0^2 - \lambda^2)$，由此得
\[
u_\pm = u_r \pm 2\lambda\sqrt{\omega_0^2 - \lambda^2},
\qquad
\Omega_+^2 - \Omega_-^2 = 4\lambda\sqrt{\omega_0^2 - \lambda^2} .
\]
于是 $\Omega_+ - \Omega_- = \frac{\Omega_+^2 -
\Omega_-^2}{\Omega_+ + \Omega_-}$，且当 $\lambda \ll \omega_0$ 时
二者均有 $\Omega_\pm \approx \omega_0$：$\Omega_+ - \Omega_- \approx
\frac{4\lambda\omega_0}{2\omega_0} = 2\lambda$，故
$\frac{\omega_0}{\Omega_+ - \Omega_-} \approx
\frac{\omega_0}{2\lambda} = Q$。测定共振峰的宽度，
即测得其品质因数。

\textbf{16.} $Q = 10$; $\Omega_r = \sqrt{1 - 2(0.05)^2} =
\sqrt{0.995} = 0.9975$; $R_{\max} = \frac1{2 \times 0.05
\sqrt{1 - 0.0025}} = \frac1{0.1 \times 0.99875} = 10.01$；静态
响应 $R(0) = 1$；带宽 $\approx 2\lambda = 0.1$。一个高而
窄的尖峰，高 $\approx Q$，立于高为 $1$ 的平台之上。

\textbf{17.} 若 $2\lambda^2 \geq \omega_0^2$，则 $u_r \leq 0$，且
$g'(u) = 2(u - u_r) > 0$ 对所有 $u > 0$ 成立，故 $g$ 在
$\intoo0{+\infty}$ 上严格递增，于是 $R = A/\sqrt g$ 自
$R(0) = A/\omega_0^2$ 起严格递减：响应在零频率处最大，
且不存在峰值。

\textbf{18.} $\gamma = \iu\Omega$ 不是 $r^2 +
\omega_0^2$ 的根（因 $\Omega \neq \omega_0$），于是
由\cref{met:b1:diffeq:particular}，取 $m = 0$，得 $x_p =
\frac{A\cos\Omega t}{\omega_0^2 - \Omega^2}$（代入检验：
$-\Omega^2 + \omega_0^2$ 乘以余弦）。通解
为：
\[
x(t) = \frac{A\cos\Omega t}{\omega_0^2 - \Omega^2}
+ \lambda_1\cos\omega_0 t + \mu_1\sin\omega_0 t .
\]

\textbf{19.} 由 $x(0) = 0$ 可得 $\lambda_1 = -\frac A{\omega_0^2 -
\Omega^2}$，又由 $x'(0) = 0$ 得 $\mu_1 = 0$：
\[
x(t) = \frac{A}{\omega_0^2 - \Omega^2}\,
\bigl(\cos\Omega t - \cos\omega_0 t\bigr)
= \frac{2A}{\omega_0^2 - \Omega^2}\,
\sin\Bigl(\frac{(\omega_0 - \Omega)t}2\Bigr)
\sin\Bigl(\frac{(\omega_0 + \Omega)t}2\Bigr) ,
\]
由和差化积公式 $\cos a - \cos b =
2\sin\frac{b - a}2\sin\frac{b + a}2$，取 $a = \Omega t$，
$b = \omega_0 t$。

\textbf{20.} 当 $\Omega$ 接近 $\omega_0$ 时，第二个正弦
以快频率 $\frac{\omega_0 + \Omega}2
\approx \omega_0$ 振动，而第一个因子是频率为
$\frac{\abs{\omega_0 - \Omega}}2$ 的慢包络，快
振动的振幅随包络周期而增减，
$\frac{2\pi}{\abs{\omega_0 - \Omega}}$（每个包络周期含两次
拍），达到极大值 $\frac{2A}{\abs{\omega_0^2 - \Omega^2}}$。
当 $\Omega \to \omega_0$ 时，拍既变慢（周期
$\to \infty$），又变高（振幅 $\to \infty$）。

\textbf{21.} 固定 $t$。当 $\Omega \to \omega_0$ 时：
\[
\frac{2A}{\omega_0^2 - \Omega^2}\sin\Bigl(\frac{(\omega_0 -
\Omega)t}2\Bigr)
= \frac{2A}{\omega_0 + \Omega}\cdot
\frac{\sin\bigl(\frac{(\omega_0 - \Omega)t}2\bigr)}
{\omega_0 - \Omega}
\longrightarrow \frac{2A}{2\omega_0}\cdot\frac t2
= \frac{At}{2\omega_0} ,
\]
而 $\sin\bigl(\frac{(\omega_0 + \Omega)t}2\bigr) \to
\sin(\omega_0 t)$，故极限为 $x_\infty(t) =
\frac{At\sin\omega_0 t}{2\omega_0}$。直接验证：取 $C =
\frac A{2\omega_0}$，$x_\infty = Ct\sin\omega_0 t$ 有
$x_\infty'' = 2C\omega_0\cos\omega_0 t - C\omega_0^2 t
\sin\omega_0 t$，于是 $x_\infty'' + \omega_0^2 x_\infty =
2C\omega_0\cos\omega_0 t = A\cos\omega_0 t$，又 $x_\infty(0) =
0$，且 $x_\infty'(0) = C\sin 0 + C\omega_0 \cdot 0 \cdot \cos 0 =
0$ --- 当 $\omega_0 = A = 1$ 时，这恰好
就是\cref{ex:b1:diffeq:oscillation}。共振是
拍的退化形式：包络的第一次隆起，
被拉伸至无穷长。

\textbf{22.} 无阻尼时，共振振幅线性增长
且无界；有阻尼 $\lambda > 0$ 时，增长
止于 $R_{\max} \approx Q\,\frac A{\omega_0^2}$。其机制为：
耗散消耗能量的速率随振幅增大而增大，
（第 23 题），故积累恰在阻尼
消耗能量的速率与外力供给能量的速率相等时停止。

\textbf{23.} 记 $x_p' = -R\Omega\sin(\Omega t - \varphi)$，在一个周期内
平均值 $\langle\cos^2\rangle =
\langle\sin^2\rangle = \frac12$，且 $\langle\sin\cos\rangle = 0$，
由此得：
\[
\langle P_{\mathrm{diss}}\rangle
= 2\lambda\,R^2\Omega^2\,\langle\sin^2(\Omega t -
\varphi)\rangle = \lambda R^2\Omega^2 ;
\]
再将 $\sin(\Omega t - \varphi) = \sin\Omega t\cos\varphi
- \cos\Omega t\sin\varphi$：
\[
\langle P_{\mathrm{in}}\rangle
= -AR\Omega\,\bigl\langle\cos\Omega t\,\sin(\Omega t -
\varphi)\bigr\rangle
= AR\Omega\,\frac{\sin\varphi}2 .
\]
由于 $\sin\varphi = \frac{2\lambda\Omega}{\sqrt D} =
\frac{2\lambda\Omega R}A$，于是 $\frac{AR\Omega}2 \cdot
\frac{2\lambda\Omega R}A = \lambda R^2\Omega^2$：输入与
耗散功率恰好相等 --- 这正是一个
稳态所具有的定义性质。

\textbf{24.} （i）结构定理将每个解分解为
稳态与瞬态之和（第 9--11 题），并将
唯一性归结为齐次问题。（ii）复数方法
把求特解的过程化为一次复数除法
（第 6 题），振幅与相位由
模与辐角读出。（iii）共振曲线是一个纯粹的
函数研究 --- 关于 $u = \Omega^2$ 的二次式、其极小值、
其水平集 --- 沿用\cref{ch:b1:functions}中的
方法（第 12--17 题）。

\textbf{25.} 自由且有阻尼：衰减的伪振荡，
寿命为 $\frac1\lambda$，约 $\frac Q\pi$ 次摆动。受迫且
有阻尼：暂态消亡，唯一的正弦稳态
在外力频率下保留，振幅在
$\omega_0$ 附近达到峰值（高度 $\approx Q \times$ 静态值，宽度 $\approx
\frac{\omega_0}Q$），相位自 $0$ 变至 $\pi$，并经过
$\frac\pi2$，恰在 $\omega_0$ 处。自由且无阻尼：永续
振荡。受迫且无阻尼：拍，退化为
精确调谐时的线性增长共振。一个无量纲
数 $Q = \frac{\omega_0}{2\lambda}$ 调控全部 ---
峰高、带宽与暂态寿命是同一旋钮上的三种
读数。后续：将 $x'' + 2\lambda x' +
\omega_0^2x$ 改写为一阶方程组，即可引出
\cref{ch:b1:matrices}中的矩阵方法及本科二年级分册；而将一个任意
周期外力分解为正弦之和（傅里叶级数，本科三年级分册），
使本题的单频分析成为普遍的基本构件
：对每个频率求解，再叠加。
\end{solution}
